%%
%% Copyright 2019-2020 Elsevier Ltd
%%
%% This file is part of the 'CAS Bundle'.
%% --------------------------------------
%%
%% It may be distributed under the conditions of the LaTeX Project Public
%% License, either version 1.2 of this license or (at your option) any
%% later version.  The latest version of this license is in
%%    http://www.latex-project.org/lppl.txt
%% and version 1.2 or later is part of all distributions of LaTeX
%% version 1999/12/01 or later.
%%
%% The list of all files belonging to the 'CAS Bundle' is
%% given in the file `manifest.txt'.
%%
%% Template article for cas-sc documentclass for
%% double column output.

%\documentclass[a4paper,fleqn,longmktitle]{cas-sc}
\documentclass[a4paper,fleqn]{cas-sc}
\usepackage{njfutitlepage}
% \usepackage[numbers]{natbib}
\usepackage[authoryear]{natbib}
% \usepackage[authoryear,longnamesfirst]{natbib}
\usepackage{csquotes}
\usepackage[UTF8]{ctex}
\usepackage{gbt7714}
\usepackage{makecell}

\begin{document}
\let\WriteBookmarks\relax
\def\floatpagepagefraction{1}
\def\textpagefraction{.001}

\shortauthors{7组（王孜悠）}
\shorttitle{}
\title [mode = title]{\makecell[c]{FeHALS：基于Web的3D可视化航路规划与激光仿真系统\\\textbf{技术设计书}}}

\Course{GIS设计与开发}
\Major{测绘工程}
\Grade{2023级}
\Group{7 组}

\Member[2310604116]{王 孜悠}
\Member[2310604109]{梅 康凯}
\Member[2310604115]{唐 如意}
\Member[2310604122]{杨 杰瑞}
\Member[2310604117]{温 璞}
\Member[2110604112]{刘 继明}
\Member[2110604106]{李 英杰}
\Member[2310604103]{杜 仲宇}
\Member[2310604113]{施 祺}
\Member[2310604118]{向 宇鸣}

\Teacher{隋 铭明}
\Teacher{那 嘉明}
\Teacher{陈 动}
\Teacher{朱 杰}

\MakeNJFUTitlePage[UseDocTitle]
\tableofcontents
\newpage

\section{背景目的}
1.1. 项目背景
机载激光扫描（Airborne Laser Scanning，ALS）能够直接获取地表及地物的三维离散采样，是数字
高程模型生产、森林调查、城市三维建模和灾害监测的重要数据源。真实飞行采集受设备成本、空域、天
气和任务风险制约，外业前通过虚拟场景试验航线与扫描参数，可以降低反复试飞的成本。HELIOS++
是面向地面、移动、无人机及机载激光扫描的开源仿真框架，支持场景、平台、扫描器和测量任务的组
合配置 [1-2]，但其命令行和 XML 配置方式对初学者存在较高门槛。
FeHALS（Frontend of HELIOS++ for Airborne Laser Scanning）以浏览器为操作入口，把三维预
览、航点编辑、弓字形航线生成、参数配置、HELIOS++ 任务执行、实时日志和点云分析串联为完整
流程。开放式 Web GIS 技术便于跨平台访问与后续扩展 [3-5]，浏览器侧利用 WebGL 和 Three.js 完
成硬件加速三维显示，无需安装专用桌面 GIS 软件。
1.2. 建设目标
本项目旨在为 HELIOS++ 提供直观、可交互、可追踪的前端工作台：集中管理模型、航迹和扫
描参数；在 Z-up 三维场景中完成航点规划；生成 HELIOS++ 所需轨迹和 XML；以异步任务运行仿
真并实时显示日志；解析 XYZ、LAS、LAZ 结果，提供点云渲染与统计；形成便于多人协作和迭代的
Web GIS 工程结构。
系统面向激光雷达仿真学习者、测绘工程学生和方案验证人员，定位为教学与原型验证工具，不
替代实际航空作业中的空域审批、飞控、坐标转换及成果质量验收。本设计书以 GitHub 仓库 master
分支当前实现为准，规划中的 LOD、任务队列、Docker 等功能不会当作现成功能描述。
1.3. 设计原则
系统遵循坐标与参数一致、操作可视化、任务可追踪、前后端模块化和能力渐进增强原则。前端、
后端和 HELIOS++ 均以 Z 轴表示高程；任务具有唯一标识、状态、日志和结果；组件、状态、接口
和服务分层；当前先满足课程演示和百万点级预览，超大点云与生产部署作为后续建设内容。

\section{需求分析}
2. 需求分析
2.1. 业务流程
一次完整业务流程为：启动前后端并检查 HELIOS++ 环境；上传场景模型并检查方向和包围盒；
手工编辑航点或由矩形生成弓字形航线；设置平台、飞行与扫描参数；生成轨迹和配置；提交仿真并查
看实时日志；加载结果点云、查看统计与直方图；下载结果或清理缓存。
2.2. 功能需求
\begin{table}[htbp]
  \centering
  \caption{FeHALS功能需求}
  \begin{tabular}{|c|l|}
    \hline
    \textbf{编号} & \textbf{需求说明} \\
    \hline
    FR-01 & 上传OBJ、GLTF、GLB、STL；显示/隐藏/删除模型，显示包围盒。仅OBJ作为当前HELIOS++场景输入。 \\
    FR-02 & 三维场景旋转、缩放、平移、视角适配、Z-up网格和模型方向转换。 \\
    FR-03 & 支持点击添加、拖动、右键或Delete删除、列表编辑和清空，并以折线连接。 \\
    FR-04 & 选择矩形范围和航线间距，自动生成连续往返的弓字形航点。 \\
    FR-05 & 以所有模型世界包围盒的最大Z加20m计算建议高度，并受平台范围约束。 \\
    FR-06 & 设置UAV/Airborne、速度、高度、扫描频率、扫描半角、脉冲频率和输出格式。 \\
    FR-07 & 生成恒定高度轨迹文件和HELIOS++ survey XML。 \\
    FR-08 & 异步调用HELIOS++，记录pending/running/completed/failed状态并支持取消。 \\
    FR-09 & WebSocket实时推送日志，REST接口可补取历史日志。 \\
    FR-10 & 查询、下载和解析结果；显示点数、边界、高程/强度统计及40组直方图。 \\
    FR-11 & 统计并分类清理models, trajectories, configs, results。 \\
    \hline
  \end{tabular}
  \label{tab:functions}
\end{table}
2.3. 非功能需求与约束
易用性要求核心流程在单页完成，参数显示单位、范围和默认值；性能要求全量统计与显示数据
分离，前端渲染最多一百万点，XYZ 每五十万行分块读取；可靠性要求子进程失败不拖垮 API，模型
异步加载期间删除不产生残留对象；可维护性要求视图、Pinia 状态、API 客户端、路由、数据模型和
业务服务解耦。
系统依赖正确安装的 HELIOS++ 及其环境路径，未配置时只能使用三维预览。UAV 与 Airborne
当前均引用 copter_linearpath 平台和 riegl_vux-1uav 扫描器，主要由参数范围区分。航点在 𝑧 = 0
地面编辑，生成轨迹时统一叠加飞行高度，属于恒高简化模型。仓库说明记录本机构建的 LASlib 输出
异常，因此默认 XYZ；LAS/LAZ 只有在 HELIOS++ 正确链接 LASlib 时可用。任务和配置主要存于
单进程内存，系统也未设置账户权限，生产部署前必须补充持久化、鉴权、上传限制、HTTPS 和审计。

\section{总体设计}
3.1. 系统架构
FeHALS 采用 B/S 三层架构：Vue 3、Three.js、Pinia 和 Axios 构成表现层；FastAPI、Pydantic、
asyncio 与 WebSocket 构成服务层；HELIOS++ 及其平台、扫描器、场景资产构成计算层。本地目录
保存模型、轨迹、配置和结果，模型 manifest 用于服务重启后的恢复。前端只交换 JSON、文件和任
务标识，后端负责文件组织、配置生成、进程管理和结果解析。
\begin{figure}[htbp]
  \centering
  \caption{FeHALS总体架构}
  \includegraphics[width=0.85\textwidth]{figs/gernal_fehals_arch.pdf}
  \label{fig:arch}
\end{figure}
前端中 Scene3D.vue 管理画布，ControlPanel.vue 管理参数，WaypointList.vue 管理航点，
LogConsole.vue 展示日志，PointCloudPanel.vue 展示结果。共享状态由 scene、waypoints、simulation
三个 store 管理；跨组件逻辑放在 composables。后端按 API 路由、Pydantic schema、业务 service 和
config 分层。
3.2. 数据与控制流程
组件操作先更新 Pinia，再由 useHeliosAPI 请求后端。后端为轨迹、配置和任务分配 UUID，生
成文件并启动 HELIOS++ 子进程；日志写入任务缓冲并广播；成功后定位结果文件、计算全量统计并
返回显示抽样；前端创建 THREE.Points 绘制点云。
开发环境中 Vite 默认使用 5173 端口，并把 /api 和 /ws 转发到 8000 端口 FastAPI；后端通过
/static 发布文件。Windows 用 run.bat，类 Unix 用 run.sh。当前是单机双进程课程演示结构，尚
未包括数据库、反向代理、容器编排和集中监控。
\begin{figure}[htbp]
  \centering
  \caption{FeHALS模块架构}
  \includegraphics[width=\textwidth]{figs/detialed_fehals_arch.pdf}
  \label{fig:module}
\end{figure}
3.3. 关键设计决策
采用 Web 前端可以降低分发门槛并把场景、参数与日志放在同一工作区 [4-5]。系统显式将相机
up 设为 (0, 0, 1)，Y-up 模型加载时绕 X 轴旋转 90◦，保证测绘高程语义一致。资源操作采用 REST，
连续日志采用 WebSocket[6-7]。GLTF/GLB/STL 只承担当下前端预览，避免与 HELIOS++ 输入能力
混淆。点云采取“全量统计、抽样渲染”，在统计准确性与交互性能间折中；更大数据仍需要层次 LOD
或渐进渲染 [8-9]。
\begin{figure}[htbp]
  \centering
  \caption{FeHALS工作流程}
  \includegraphics[width=0.85\textwidth]{figs/gernal_fehals_workflow.pdf}
  \label{fig:workflow}
\end{figure}
\section{功能设计}
4.1. 模型、场景与航点
后端校验扩展名并以随机标识保存上传文件，返回静态 URL。前端按扩展名选择 OBJLoader、
GLTFLoader 或 STLLoader，将对象加入 modelsGroup。模型可见性和包围盒可切换，Box3 计算尺
寸。若模型在异步加载完成前已被删除，pendingRemoval 会使迟到结果被释放，避免“幽灵对象”。相
机按包围盒自适应，地面为 XY 平面，Z 为高程。
Raycaster 与 𝑧 = 0 隐形地面求交：点击添加航点，拖动修改 XY，右键或 Delete 删除，红色球体
和折线表示路径，选中航点变黄。拖动时临时禁用相机控制，以 25 像素平方阈值区分点击和拖动。弓
字形航线在矩形内按间距 𝑑 生成，条数
N = \left| \frac{y_{\max} - y_{\min}}{d} \right| + 1（1）
相邻扫描线的 X 端点次序交替，形成连续往返路径。
4.2. 参数、轨迹和配置
建议高度取场景包围盒并集最高点 𝑧𝑚𝑎𝑥 加 20 m：
h_{\text{suggest}} = z_{\max} + 20\text{ m} (2)
UAV 速度为 0.5–50 m/s、高度为 3–500 m；Airborne 速度为 10–200 m/s、高度为 3–5000 m。RIEGL
VUX-1UAV 扫描频率为 10–200 Hz，扫描半角为 1
◦–165◦，脉冲频率为 50–550 kHz。后端将航点改写为统一高度的轨迹文件，并把平台、扫描器、轨迹、速度、频率、半角和格式写入 survey XML。轨迹
生成与配置生成分离，便于复用和排错。
4.3. 仿真、日志和结果
任务提交后检查 trajectory、config 和可选 scene model，分配 UUID 并由 pending 进入 running。
HELIOS++ 标准输出和错误输出逐行保存并广播到 /ws/logs/\{task_id\}；退出码为 0 且找到结果
时进入 completed，否则进入 failed。运行中任务可取消。前端保存 taskId、状态、进度、日志和结果，
WebSocket 不可用时可通过 REST 查询。
LAS/LAZ 由 laspy 读取，XYZ 分块读取并过滤非法记录。统计包括总点数、XYZ 边界、高程均
值、标准差、最小/最大、中位数、P5/P95、40 组高程直方图，以及存在强度字段时的强度统计。渲
染超过一百万点时使用覆盖全文件序列的等间隔索引抽样。前端以 BufferGeometry 和 PointsMaterial
绘制，支持固定色、高度色和强度色；归一化为
t = \frac{v - v_{\min}}{\max(v_{\max} - v_{\min}, 1)}, (3)
再映射至蓝—青—绿—黄—红色带。抽样仅影响显示，不改变全量统计。
4.4. 接口设计
\begin{table}[htbp]
  \centering
  \caption{FeHALS后端接口}
  \begin{tabular}{|c|l|l|}
    \hline
    \textbf{方法} & \textbf{路径} & \textbf{说明} \\
    \hline
    POST & /api/models/upload & 上传模型 \\
    GET  & /api/models & 查询模型列表 \\
    DELETE & /api/models/\{model\_id\} & 删除模型 \\
    POST & /api/trajectory/generate & 生成航迹 \\
    GET  & /api/trajectories & 查询航迹列表 \\
    GET  & /api/trajectories/download/\{file\_id\} & 下载航迹 \\
    POST & /api/config/generate & 生成配置 \\
    POST & /api/simulation/run & 创建仿真任务 \\
    GET  & /api/simulation/status/\{task\_id\} & 查询任务状态 \\
    GET  & /api/simulation/logs/\{task\_id\} & 获取任务日志 \\
    POST & /api/simulation/cancel/\{task\_id\} & 取消任务 \\
    GET  & /api/results/\{task\_id\} & 获取统计和抽样 \\
    GET  & /api/results/\{task\_id\}/download & 下载结果 \\
    GET  & /api/cache & 查看缓存统计 \\
    DELETE & /api/cache/\{cache\_type\} & 分类清理缓存 \\
    \hline
  \end{tabular}
  \label{tab:apis}
\end{table}
\begin{table}[htbp]
  \centering
  \caption{主要数据对象}
  \begin{tabular}{|c|c|c|c|}
    \hline
    \textbf{对象} & \textbf{形式} & \textbf{关键字段} & \textbf{用途} \\
    \hline
    场景模型 & OBJ/GLTF/GLB/STL & id, name, url, up, bbox & 前端均可预览，当前只有OBJ进入仿真。 \\
    航点 & JSON & x, y, z & Pinia中编辑，生成轨迹前可反复修改。 \\
    轨迹 & 文本 & id, altitude, path & 由航点生成，供配置引用。 \\
    配置 & XML & platform, scanner, speed, angle, frequency & 供HELIOS++读取。 \\
    任务 & 内存对象 & task\_id, status, logs, result & 当前随服务进程生命周期存在。 \\
    结果 & XYZ/LAS/LAZ & X, Y, Z, 可选intensity & 保存原始文件并产生统计和显示抽样。 \\
    \hline
  \end{tabular}
  \label{tab:dataobjects}
\end{table}
后端用 Pydantic 进行结构校验，资源不存在返回 404，参数或格式错误返回 4xx，子进程失败保
留日志。前端在上传、加载、请求和连接处捕获异常并给出可操作提示。公网部署时不应暴露服务器绝
对路径和内部堆栈。
\section{数据分析}
5.1. 数据对象与组织
后端请求模型包括 TrajectoryRequest、ConfigRequest、SimulationRunRequest 和 SimulationStatus。前端 scene store 管理模型和点云选项，waypoints store 管理航点，simulation store 管理任务、日
志、结果与参数。大型 Three.js 对象由单例 useThreeScene 管理，store 只保存业务镜像，避免响应
式系统深度代理图形对象。
5.2. 坐标、统计与性能
距离单位为米，速度为 m/s，扫描频率为 Hz，脉冲频率为 kHz，扫描角为半角度数。系统当前不
读取 EPSG，也不进行投影转换，因此模型和轨迹必须预先处于同一以米为单位的坐标系；后续应增
加 CRS 元数据和 pyproj 转换。设点数为 𝑛、高程为 𝑧𝑖，则
\bar{z} = \frac{1}{n}\sum_{i=1}^{n} z_i, \qquad 
\sigma_z = \sqrt{\frac{1}{n}\sum_{i=1}^{n} (z_i - \bar{z})^2}(4)
中位数、P5 和 P95 描述主体分布并降低极端值影响。若 𝑛 > 𝑚 且 𝑚 = 1,000,000，显示索引可取
i_k = \left\lfloor \frac{k(n-1)}{m-1} \right\rfloor, \quad k = 0, \dots, m-1 (5)
该方法简单确定，但并非空间均匀采样；超大点云仍需八叉树、多分辨率或渐进策略 [8-9]。
后端按 models、trajectories、configs、results 建目录。数据依赖为“场景和航点 → 轨迹与参数 →
配置 → 任务 → 点云”。生产版本应记录资源引用，避免清理活动任务依赖。上传侧还需限制大小、校
验 MIME 和解析结果、防路径穿越；结果下载只能通过任务记录定位。质量检查至少覆盖航点数、有
限坐标、飞行高度、扫描器范围、文件存在、退出码和非空结果。

\section{主要技术手段}
6.1. 前端与三维显示
Vue 3 组合式 API 构建单页应用，Vite 提供开发与构建，Pinia 管理共享状态，Axios 统一请求。
Three.js 封装相机、光照、材质、Loader、BufferGeometry、Raycaster 和 OrbitControls，实现模型、
航点和点云交互。点云使用单一缓冲几何，避免逐点对象开销。WebGL 无插件硬件加速适合 Web 端
地形点云显示 [9]，后续可将解析放入 Web Worker，并引入八叉树 LOD 或渐进渲染 [8]。
6.2. 后端与 HELIOS++
FastAPI 提供异步路由与 OpenAPI 文档，Pydantic 在业务执行前完成类型校验，asyncio 子进程
避免长任务阻塞 API，WebSocket 负责连续日志。REST 有利于资源接口解耦 [6]，WebSocket 适合持
续消息 [7]。HELIOS++ 使用 C++ 实现，支持多类场景、平台和扫描器，并允许在物理真实性和计
算成本间配置 [2]。本系统采用命令行子进程集成，便于故障隔离和日志记录；平台/扫描器范围来自
HELIOS++ 资产 XML。
LAS/LAZ 由 laspy 解析，XYZ 以文本流读取，NumPy 计算统计。Web 端 LiDAR 系统应明确处
理、存储和分发边界 [5]；开源 Web 架构可将后端处理与浏览器可视分析组成完整流程 [4]。
6.3. 协作、测试与改进
仓库通过 GitHub 分支、Pull Request 和审查协作；README 说明安装，CHANGELOG 记录演
进，SUMMERY 总结阶段成果，TODO 区分已实现与计划项。交付前应补充：轨迹、弓字形、参数边界、
XML 转义、XYZ 解析和统计单元测试；所有 REST 路由的接口测试；小型 OBJ 场景的 HELIOS++
集成测试；模型增删、航点拖拽、日志重连、点云着色的前端测试；10 万、100 万及更大点云的耗时、
帧率和内存测试。
后续应优先完成 HELIOS++ 环境诊断、统一错误结构与自动测试；再将任务、配置和资源关系持
久化并引入任务队列；随后增加 CRS、空间索引、LOD、覆盖热力图和仿真动画；最后完成 Docker、
反向代理、鉴权、HTTPS、配额与监控。当前 FeHALS 已打通“可视化规划—配置—仿真—日志—点
云分析”的核心链路，适合作为课程设计和 HELIOS++ 入门工作台。

\newpage
\bibliographystyle{gbt7714-numerical}
\bibliography{cas-refs}

\end{document}
